\hypertarget{c11-core-language-features}{%
\chapter{C++11 CORE LANGUAGE
FEATURES}\label{c11-core-language-features}}

The C++11 standard (originally known as C++0x) was a revolutionary
update that transformed C++ into a modern language. It addressed
verbosity, safety, and performance.

\begin{center}\rule{0.5\linewidth}{0.5pt}\end{center}

\hypertarget{auto-type-deduction}{%
\section{1. AUTO \& TYPE DEDUCTION}\label{auto-type-deduction}}

\hypertarget{the-auto-keyword}{%
\subsection{\texorpdfstring{1.1 The \texttt{auto}
Keyword}{1.1 The auto Keyword}}\label{the-auto-keyword}}

In C++98, you had to explicitly declare types, which could be verbose,
especially with iterators. C++11 introduced
\passthrough{\lstinline!auto!} to ask the compiler to deduce the type
from the initializer.

\begin{lstlisting}[language={C++}]
// C++98
int x = 5;
std::vector<int> v;
std::vector<int>::iterator it = v.begin();

// C++11
auto x = 5;       // deduced as int
auto it = v.begin(); // deduced as std::vector<int>::iterator
\end{lstlisting}

\textbf{Key Rules:} 1. \textbf{Must be initialized}:
\passthrough{\lstinline!auto x;!} is invalid. 2. \textbf{References}:
\passthrough{\lstinline!auto!} strips references. Use
\passthrough{\lstinline!auto\&!} to keep them. 3. \textbf{Const}:
\passthrough{\lstinline!auto!} strips top-level const. Use
\passthrough{\lstinline!const auto!} or
\passthrough{\lstinline!const auto\&!}.

\begin{lstlisting}[language={C++}]
int a = 10;
int& ref = a;

auto x = ref;  // x is int (copy), not int&
auto& y = ref; // y is int& (reference)
\end{lstlisting}

\hypertarget{decltype}{%
\subsection{\texorpdfstring{1.2
\texttt{decltype}}{1.2 decltype}}\label{decltype}}

\passthrough{\lstinline!decltype!} inspects the declared type of an
entity or expression. Unlike \passthrough{\lstinline!auto!}, it does not
strip references or const.

\begin{lstlisting}[language={C++}]
int x = 5;
decltype(x) y = 10; // y is int

const int& z = x;
decltype(z) w = x; // w is const int&
\end{lstlisting}

\begin{center}\rule{0.5\linewidth}{0.5pt}\end{center}

\hypertarget{range-based-for-loops}{%
\section{2. RANGE-BASED FOR LOOPS}\label{range-based-for-loops}}

C++11 introduced a syntactic sugar for iterating over containers,
arrays, and initializer lists.

\begin{lstlisting}[language={C++}]
std::vector<int> vec = {1, 2, 3, 4, 5};

// C++98
for (std::vector<int>::iterator it = vec.begin(); it != vec.end(); ++it) {
    std::cout << *it << " ";
}

// C++11 Range-based for
for (int i : vec) {
    std::cout << i << " ";
}

// With auto (Common pattern)
for (auto i : vec) {
    std::cout << i << " ";
}

// By Reference (to modify elements)
for (auto& i : vec) {
    i *= 2;
}

// By Const Reference (read-only, avoids copy)
for (const auto& i : vec) {
    std::cout << i << " ";
}
\end{lstlisting}

\begin{center}\rule{0.5\linewidth}{0.5pt}\end{center}

\hypertarget{uniform-initialization}{%
\section{3. UNIFORM INITIALIZATION}\label{uniform-initialization}}

C++11 introduced a consistent syntax for initializing everything using
curly braces \passthrough{\lstinline!\{\}!}.

\hypertarget{brace-initialization}{%
\subsection{3.1 Brace Initialization}\label{brace-initialization}}

\begin{lstlisting}[language={C++}]
// C++98 inconsistency
int a = 5;
int arr[] = {1, 2};
std::vector<int> v; v.push_back(1); // Tedious

// C++11 Uniformity
int a{5};
int arr[]{1, 2};
std::vector<int> v{1, 2, 3}; // Initializer list!
std::string s{"Hello"};
\end{lstlisting}

\hypertarget{preventing-narrowing}{%
\subsection{3.2 Preventing Narrowing}\label{preventing-narrowing}}

Brace initialization forbids ``narrowing conversions'' where data loss
might occur.

\begin{lstlisting}[language={C++}]
int x = 3.14;  // C++98: Compiles (x becomes 3), implicit cast
// int y{3.14}; // C++11: Error! Narrowing conversion
\end{lstlisting}

\hypertarget{stdinitializer_list}{%
\subsection{\texorpdfstring{3.3
\texttt{std::initializer\_list}}{3.3 std::initializer\_list}}\label{stdinitializer_list}}

Classes can now take a list of elements in their constructor.

\begin{lstlisting}[language={C++}]
#include <initializer_list>

class MyClass {
public:
    MyClass(std::initializer_list<int> list) {
        for (auto elem : list) {
            // process elem
        }
    }
};

MyClass obj = {1, 2, 3, 4, 5};
\end{lstlisting}

\begin{center}\rule{0.5\linewidth}{0.5pt}\end{center}

\hypertarget{nullptr}{%
\section{4. NULLPTR}\label{nullptr}}

C++98 used \passthrough{\lstinline!0!} or \passthrough{\lstinline!NULL!}
(macro for 0) for null pointers. This caused ambiguity with function
overloading.

\begin{lstlisting}[language={C++}]
void func(int x) { std::cout << "Integer"; }
void func(char* p) { std::cout << "Pointer"; }

// C++98
func(NULL); // Calls func(int)! Because NULL is 0.

// C++11
func(nullptr); // Calls func(char*).
\end{lstlisting}

\passthrough{\lstinline!nullptr!} is a keyword of type
\passthrough{\lstinline!std::nullptr\_t!}. It implicitly converts to any
pointer type but \emph{not} to integers (except bool
\passthrough{\lstinline!false!}).

\begin{center}\rule{0.5\linewidth}{0.5pt}\end{center}

\hypertarget{strongly-typed-enums}{%
\section{5. STRONGLY TYPED ENUMS}\label{strongly-typed-enums}}

C++98 enums leaked their names into the surrounding scope and implicitly
converted to integers. C++11 \passthrough{\lstinline!enum class!} fixes
this.

\begin{lstlisting}[language={C++}]
// C++98
enum Color { RED, GREEN, BLUE };
int r = RED; // Implicit conversion allowed

// C++11
enum class TrafficLight { RED, YELLOW, GREEN };

// TrafficLight t = RED; // Error: RED not in scope
TrafficLight t = TrafficLight::RED;
// int x = TrafficLight::RED; // Error: No implicit conversion

// Explicit underlying type
enum class Byte : unsigned char { A, B, C };
\end{lstlisting}

\begin{center}\rule{0.5\linewidth}{0.5pt}\end{center}

\hypertarget{other-core-features}{%
\section{6. OTHER CORE FEATURES}\label{other-core-features}}

\begin{itemize}
\tightlist
\item
  \textbf{\passthrough{\lstinline!constexpr!}}: Allows functions and
  variables to be evaluated at compile-time (limited in C++11, expanded
  later).
\item
  \textbf{\passthrough{\lstinline!static\_assert!}}: Compile-time
  assertion checking.
  \passthrough{\lstinline!cpp     static\_assert(sizeof(int) == 4, "Int must be 4 bytes");!}
\item
  \textbf{Delegating Constructors}: One constructor can call another.
\item
  \textbf{\passthrough{\lstinline!default!} and
  \passthrough{\lstinline!delete!} functions}:
  \passthrough{\lstinline!cpp     class NonCopyable \{         NonCopyable(const NonCopyable\&) = delete; // Ban copying         NonCopyable() = default; // Explicitly request default ctor     \};!}
\item
  \textbf{\passthrough{\lstinline!override!} and
  \passthrough{\lstinline!final!}}: Virtual function controls (See
  Chapter 19).
\end{itemize}
